% Generated draft: requires mathematical and source review.
\documentclass[11pt]{article}
\usepackage[margin=25mm]{geometry}
\usepackage{fontspec}
\setmainfont{Latin Modern Roman}
\usepackage{amsmath,amssymb,mathtools,unicode-math}
\setmathfont{Latin Modern Math}
\usepackage{enumitem,needspace,xurl,hyperref}
\setlength{\parindent}{0pt}
\setlength{\parskip}{5pt}
\setlength{\emergencystretch}{3em}
\newcommand{\N}{\mathbb{N}}
\newcommand{\Z}{\mathbb{Z}}
\newcommand{\Q}{\mathbb{Q}}
\newcommand{\R}{\mathbb{R}}
\newcommand{\C}{\mathbb{C}}
\newcommand{\F}{\mathbb{F}}
\newcommand{\A}{\mathbb{A}}
\newcommand{\PP}{\mathbb{P}}
\newcommand{\E}{\mathbb{E}}
\newcommand{\eps}{\varepsilon}
\newcommand{\dd}{\,\mathrm{d}}
\begin{document}
\section*{UnsolvedMath ID: 30002660}
\textbf{Source problem code:} catalogue entry
\subsection{Definitions and mathematical statement}
A Euclidean convex body conventionally means a compact convex set \(K\subset\R^n\) with nonempty interior. The supplied record does not specify the dimension \(n\), boundary regularity, or the precise class of billiard trajectories under consideration.

To describe the intended minimization problem schematically, let \(D\) be a billiard domain in an ambient space \(M\), and let \(\mathcal B(D)\) denote a specified class of admissible billiard trajectories in \(D\). For a trajectory \(\gamma\), let \(L(\gamma)\) denote its length in the ambient geometry. A shortest trajectory in this class would be an element \(\gamma_*\in\mathcal B(D)\) satisfying
\[
L(\gamma_*)=\inf_{\gamma\in\mathcal B(D)}L(\gamma).
\]
This notation is only a framework: the source does not define \(\mathcal B(D)\), assert attainment of the infimum, or specify whether trajectories must be nonconstant, closed, or periodic.

\textbf{Source question.} Can the Bezdek--Bezdek approach to shortest billiard trajectories in Euclidean convex bodies be extended to hyperbolic or spherical geometry, or to Riemannian and Finsler manifolds?

\textbf{Unresolved formulation.} The supplied record does not state which theorem, characterization, or construction constitutes the Bezdek--Bezdek approach, or what conclusions an extension must preserve. It also does not specify the admissible domains, notion of convexity, boundary conditions, reflection law, trajectory class, or hypotheses on the ambient manifolds. In the Finsler setting, the treatment of direction-dependent length and any reversibility assumption are unspecified. Consequently, the question is retained as a research direction; the supplied information does not determine a single quantified mathematical assertion. The four indicated settings are possible extension targets, and the source does not specify whether an extension in any one setting or a result covering all of them is requested.
\subsection{Short English statement}
The question asks whether the Bezdek--Bezdek method for shortest billiard trajectories in Euclidean convex bodies extends to hyperbolic, spherical, Riemannian, or Finsler geometry. The method and the required extension are not specified sufficiently to formulate a definite theorem. The supplied open status and verification labels are curation claims requiring human review.
\subsection{Sources}
\begin{itemize}
\item Local registry: \texttt{data/unsolvedmath/problems.json}, with \texttt{catalog/problems.json} as the supplementary registry specified by the project; supplied source record for permanent UnsolvedMath ID \(30002660\), titled ``Shortest Billiards Beyond Euclidean Convex Bodies,'' source problem code \texttt{catalogue entry}, in the Oberwolfach Reports Open Problems collection. These registry files have not been independently inspected for this draft.
\item Supplied report citation: \emph{Discrete Geometry (2015)}, Oberwolfach Reports, DOI \texttt{10.4171/owr/2014/40}, pp.\ 2235--2296. Supplied source URL: \url{https://doi.org/10.4171/owr/2014/40}. The background describes Oberwolfach Reports as documenting mathematical workshops and their open-problem discussions. The record separately gives proposed year \(2014\); the relationship between that year and the citation's \(2015\) label is not explained.
\item Supplied literature source, labeled ``original Oberwolfach report'': \url{https://doi.org/10.4171/owr/2014/40}. The literature assessment dated exactly \texttt{2026-08-21} is supplied dated curation, not independently verified. It labels the problem open and says that the extracted statement contained no mathematical question beyond the words ``Open Question,'' leaving its content unclear and causing normalization to open. This assessment conflicts with the explicit question in the supplied original, statement, and clean-statement fields. The background requests independent verification of time-sensitive status and cited claims.
\item The supplied statement status is \texttt{corrected\_verified}. Its verification note states that the supplied statement explicitly identifies the mathematical question and that notation and publication prose were repaired without changing its substance. This is a supplied verification claim, not verification performed for this draft. No proposer is identified. The method is attributed in the question to Bezdek--Bezdek, but no separate citation identifying their relevant result is supplied.
\end{itemize}
\end{document}
